\documentclass[lettersize,journal]{IEEEtran}
\usepackage{amsmath,amsfonts,amssymb}
\usepackage{algorithmic}
\usepackage{algorithm}
\usepackage{array}
\usepackage[caption=false,font=normalsize,labelfont=sf,textfont=sf]{subfig}
\usepackage{textcomp}
\usepackage{url}
\usepackage{graphicx}
\usepackage{cite}
\usepackage{booktabs}
\hyphenation{op-tical net-works semi-conduc-tor IEEE-Xplore}

\begin{document}

\title{Risk-Calibrated Fallback: A Selective-Prediction Approach\\
to Minimum-Risk-Manoeuvre Triggering in Closed-Loop\\
Autonomous Driving}

\author{Hassan~Waqas~Saleem, Rahat~Zeeshan~Saleem, and Hina~Maryam
\thanks{H.~W.~Saleem (corresponding author) is with the School of
Electrical Engineering and Computer Science (SEECS), National
University of Sciences and Technology (NUST), H-12, Islamabad 44000,
Pakistan (e-mail: hsaleem.msee20seecs@seecs.edu.pk).}%
\thanks{R.~Z.~Saleem is with the NUST Business School (NBS), National
University of Sciences and Technology (NUST), H-12, Islamabad 44000,
Pakistan (e-mail: rahat.mscm21nbs@student.nust.edu.pk).}%
\thanks{H.~Maryam is with the State Bank of Pakistan, I.I. Chundrigar
Road, Karachi 74000, Pakistan (e-mail: hina.maryam@sbp.org.pk).}%
\thanks{Manuscript in preparation.}}

\markboth{IEEE Transactions on Intelligent Vehicles}%
{Saleem \MakeLowercase{\textit{et al.}}: Risk-Calibrated Fallback}

\maketitle

%=======================================================================
\begin{abstract}
Every deployed autonomous vehicle has a minimum-risk manoeuvre (MRM) —
pull over, slow down, hand control back to a driver — and when to trigger
it is almost universally a hand-tuned threshold on a confidence score.
That score is precisely the quantity known to be untrustworthy under
distribution shift, so the trigger for the safety behaviour rests on a
calibration that can silently expire exactly when the vehicle needs it
most. Prior conformal-prediction work in this space is almost entirely
open-loop: it produces prediction regions and stops, so it cannot see
that a conservative region drives an MRM trigger, and that triggering
itself has a safety cost — freezing in a live lane. We formalise MRM
triggering as a selective-prediction problem with a guarantee on the
\emph{decision}, not merely on the prediction, using Learn-then-Test to
certify a trigger threshold against a bounded miss-rate target under
exchangeability, and Conformal Risk Control as a tighter alternative when
the closed-loop loss is empirically monotone in the threshold. We
evaluate in closed loop on real, replayed driving scenarios (Argoverse~2,
nuScenes) under reactive traffic, using a fallback manoeuvre compliant
with UN Regulation No.~157's minimum-risk-manoeuvre deceleration bound
and harm labelled by the nuPlan closed-loop safety definition
(at-fault collision or time-to-collision below 0.95\,s). We report both
the missed-intervention rate the calibrated trigger is built to bound and
the unnecessary-stop rate open-loop evaluation cannot see, and we show
that the two are inseparable: any trigger design that only reports one
is reporting half the picture.
\end{abstract}

\begin{IEEEkeywords}
Autonomous driving, conformal prediction, selective prediction, minimum-risk
manoeuvre, closed-loop safety evaluation, trajectory prediction, risk
calibration, distribution shift.
\end{IEEEkeywords}

%=======================================================================
\section{Introduction}
\IEEEPARstart{A}{utonomous} vehicles that lack confidence in their own
situational understanding are expected to fall back on a minimum-risk
manoeuvre: slow to a stop, pull to the shoulder, or otherwise hand
control back to a human. The decision of \emph{when} to invoke this
behaviour is, in nearly every deployed and published system, a
hand-tuned threshold on some scalar confidence or uncertainty score
produced by a perception or prediction module. That practice has a
specific and largely unexamined weakness: the very quantity the
threshold is applied to — a learned model's uncertainty estimate — is
exactly what has been shown to degrade under the kind of distribution
shift a vehicle experiences when it enters a new city, a new weather
condition, or simply a stretch of road unlike anything in its
calibration set. A threshold tuned under one operating distribution
carries no guarantee once that distribution has moved, and the failure
is silent: the system continues to report a confidence number, and nothing
in the number itself signals that the number is now wrong.

Conformal prediction offers a principled alternative: rather than a
point confidence score, it produces a prediction region with a
distribution-free, finite-sample coverage guarantee under
exchangeability. A substantial and fast-growing literature has applied
conformal prediction to trajectory forecasting and to planning under
uncertainty in autonomous driving. Almost without exception, however,
this literature is \emph{open-loop}: it constructs a prediction region,
reports its coverage and its size, and stops. Two things are invisible
from that vantage point. First, a region's effect on a downstream
control decision — specifically, whether it is used to trigger a
fallback — is not evaluated at all; open-loop coverage and closed-loop
safety outcome are simply assumed to move together. Second, and more
consequentially, \emph{triggering the fallback is not free}. A vehicle
that stops in a live lane because its uncertainty region grew large is
not obviously safer than one that continues; an evaluation that reports
only coverage, or only collision rate under the fallback, cannot see the
trade-off between the two, and cannot see that a policy which triggers
whenever it is asked to would trivially achieve a good miss rate while
being useless.

This paper studies fallback triggering as a decision problem with a
formal guarantee attached to the decision itself, evaluated where that
guarantee actually has to hold: in closed loop, on real driving
scenarios, with a trigger threshold that has consequences for what the
vehicle does next. We make three contributions.

\begin{enumerate}
\item We formalise minimum-risk-manoeuvre triggering as
\emph{selective prediction} with a guarantee on the closed-loop miss
rate — the probability that a harmful event occurs without the fallback
having fired first — rather than on marginal coverage of a static
prediction region. We calibrate the trigger threshold with
Learn-then-Test~\cite{ltt2021}, which requires no assumption that the
closed-loop loss be monotone in the threshold, and report Conformal
Risk Control~\cite{crc2022} as a tighter alternative on the
scenarios where monotonicity is observed to hold.
\item We evaluate entirely in closed loop, on scenarios replayed from
real driving logs (Argoverse~2 and nuScenes, via
ScenarioNet~\cite{scenarionet2023} and MetaDrive) under reactive traffic,
and we report the unnecessary-stop rate alongside the missed-intervention
rate at every operating point — the over-conservatism trade-off that an
open-loop study structurally cannot see.
\item We connect the closed-loop outcome to the underlying predictor's
calibration under dataset shift, asking directly whether a
degradation in open-loop conformal coverage translates into a
measurable degradation in closed-loop safety, and report the answer
honestly whether or not it supports the premise.
\end{enumerate}

The remainder of the paper is organised as follows.
Section~\ref{sec:related} situates this work against open-loop
conformal prediction for driving, task-relevant failure detection, and
closed-loop safety evaluation more broadly. Section~\ref{sec:method}
formalises the triggering problem and the calibration procedure.
Section~\ref{sec:setup} describes the closed-loop evaluation
environment, the fallback manoeuvre, and the harm labelling used to
score every rollout. Section~\ref{sec:results} reports results.
Section~\ref{sec:discussion} discusses limitations, and
Section~\ref{sec:conclusion} concludes.

%=======================================================================
\section{Related Work}
\label{sec:related}

\subsection{Conformal prediction for autonomous driving}
Conformal prediction~\cite{vovk2005} constructs prediction sets with a
distribution-free, finite-sample marginal coverage guarantee under
exchangeability between calibration and test data, and has been applied
extensively to trajectory forecasting and to safe planning under
uncertainty. Lindemann~\emph{et al.} embed conformal prediction regions
directly into a model predictive controller to obtain probabilistic
safety guarantees, but their guarantee collapses once the deployment
distribution shifts away from the calibration distribution — a
limitation Rahaman~\emph{et al.}~\cite{robustcp2026} address with robust
conformal prediction driven by a conditional generative model of the
shift, evaluated in a synthetic multi-agent corridor. SafePath
\cite{safepath2025} integrates conformal prediction into an LLM-based
path-planning pipeline, providing a guarantee that a set of candidate
paths contains a safe one with a user-specified probability, and reports
both open-loop (nuScenes) and closed-loop (highway-env, synthetic
traffic) results. In all three cases the guarantee is on the
\emph{prediction} — a region contains the truth, or a path set contains
a safe path — rather than on the eventual control decision the region is
used to make, and none evaluates a fallback-triggering decision on
real, replayed driving logs under reactive traffic. This paper's
guarantee is instead on the decision: the probability that a
minimum-risk manoeuvre fires before a harmful event, evaluated
end-to-end in closed loop.

\subsection{Task-relevant and closed-loop failure detection}
Farid~\emph{et al.}~\cite{taskrelevant2022} propose a run-time monitor
for trajectory-prediction failures that is \emph{task-relevant} by
construction: it flags a misprediction only when the resulting planning
cost is anomalous relative to the predicted cost distribution, with
data-free false-positive and false-negative rate bounds. Their detector
is evaluated against hand-labelled scenarios in an open-loop and
semi-reactive setting, and — like the conformal-prediction literature
above — does not close the loop through an actuated fallback and
re-measure the outcome. Their central insight, that a
task-irrelevant misprediction should not count against a detector, is
reflected in this paper's harm labelling: a collision the ego vehicle
suffers while already stopped or while being struck from behind is
scored as a distinct \emph{induced-collision} cost of the fallback
itself, never as a missed intervention.

\subsection{Risk-sensitive safety analysis}
Chapman~\emph{et al.}~\cite{cvar2022} develop risk-sensitive safe sets
using Conditional Value-at-Risk to assess the severity, not merely the
probability, of constraint violation along a stochastic trajectory,
providing the risk-sensitive vocabulary this paper borrows in
motivating why a mean missed-intervention rate is the appropriate
quantity to control at a chosen operating point, even though a
tail-severity extension (Section~\ref{sec:discussion}) is left to future
work.

\subsection{Closed-loop simulation infrastructure}
ScenarioNet~\cite{scenarionet2023} unifies heterogeneous real-world
driving logs — including Argoverse~2, nuScenes, nuPlan, and Lyft
Level~5 — into a common scenario description that can be replayed
interactively in the MetaDrive simulator, with either logged-trajectory
replay or reactive rule-based (IDM) traffic. This work is built directly
on that infrastructure: every scenario reported in
Section~\ref{sec:results} is a real logged driving scenario, replayed
headless with reactive traffic, not a synthetic or hand-crafted one.

%=======================================================================
\section{Problem Formulation and Method}
\label{sec:method}

\subsection{Fallback triggering as selective prediction}
\label{sec:formulation}
Consider a closed-loop rollout of a scenario $X$ under a fallback
policy governed by a scalar threshold $\lambda$: at every decision tick
the vehicle computes an uncertainty score $u_t(X)$ from a trajectory
predictor's output for nearby agents, and engages the minimum-risk
manoeuvre the first time $u_t(X) > \lambda$, remaining engaged for the
rest of the episode. Define the closed-loop loss
\begin{equation}
L(X,\lambda) = \mathbb{1}\bigl[\, \text{harm in } X \text{ under }
\lambda,\ \text{MRM not first} \,\bigr],
\label{eq:loss}
\end{equation}
the \emph{missed-intervention} indicator, and the complementary
over-conservatism quantity
\begin{equation}
S(X,\lambda) = \mathbb{1}\bigl[\, \text{MRM fires under } \lambda,\
X \text{ harm-free at } \lambda=\infty \,\bigr],
\label{eq:stop}
\end{equation}
the \emph{unnecessary stop} indicator, evaluated against the
counterfactual rollout of the same scenario with the trigger disabled.
The goal is to choose $\hat\lambda$ such that
\begin{equation}
\Pr_{X_1,\dots,X_n}\Bigl(\, \mathbb{E}_X\bigl[L(X,\hat\lambda)\bigr] \le
\alpha \,\Bigr) \ge 1-\delta,
\label{eq:guarantee}
\end{equation}
while minimising $\mathbb{E}_X[S(X,\hat\lambda)]$ — the smallest
autonomy loss compatible with the risk target $\alpha$.

Equation~\eqref{eq:guarantee} is not a conformal coverage statement:
$L(X,\lambda)$ depends on the entire closed-loop rollout under
$\lambda$, since engaging the fallback changes every subsequent state,
so it is neither a marginal coverage event on a static prediction nor
guaranteed to be monotone in $\lambda$ a priori. We therefore calibrate
$\hat\lambda$ with Learn-then-Test~\cite{ltt2021}: for a
finite grid $\lambda_1 > \lambda_2 > \cdots > \lambda_J$ ordered from
least to most conservative, compute the empirical risk
$\hat R(\lambda_j) = \tfrac{1}{n}\sum_i L(X_i,\lambda_j)$ on $n$
calibration scenarios, a Hoeffding--Bentkus $p$-value for the null
hypothesis $H_j: \mathbb{E}[L(\cdot,\lambda_j)] > \alpha$, and apply
fixed-sequence testing starting from the most conservative $\lambda$;
every $\lambda_j$ whose null is rejected satisfies
\eqref{eq:guarantee} simultaneously, with family-wise error controlled
at $\delta$, and $\hat\lambda$ is the least conservative such $\lambda_j$
(largest autonomy). Where the empirical loss is observed to be monotone
in $\lambda$ across the calibration set, we additionally report
Conformal Risk Control~\cite{crc2022}, whose expectation guarantee is
tighter but requires that monotonicity assumption; where it is violated
we report the violation rate directly rather than silently falling back
to the weaker method.

\subsection{Trigger score}
\label{sec:trigger-score}
At each decision tick the trigger score $u_t$ measures how far a
predicted future position of any nearby agent intrudes into the ego
vehicle's own near-term extent. Writing the ego's constant-heading,
constant-speed extrapolation over a horizon of $T$ steps as
$\{p^{\mathrm{ego}}_\tau\}_{\tau=1}^{T}$ — the same route-agnostic
motion model used by the harm metric in
Section~\ref{sec:harm}, deliberately \emph{not} the vehicle's planned,
route-following path, so that the trigger cannot credit itself for
steering behaviour the safety label does not assume — and a
predictor's $K$-mode forecast for agent $i$ as
$\{\hat p^{(i,k)}_\tau\}$, each predicted offset is decomposed into a
longitudinal component $a^{(i,k)}_\tau$ (along the ego heading) and a
lateral component $\ell^{(i,k)}_\tau$ (perpendicular to it), each
compared against a separate margin: a length-based longitudinal margin
$m^{\mathrm{lon}}_i = \tfrac{1}{2}(L_{\mathrm{ego}}+L_i)+c_{\mathrm{lon}}$
and a width-based lateral margin
$m^{\mathrm{lat}}_i = \tfrac{1}{2}(W_{\mathrm{ego}}+W_i)+c_{\mathrm{lat}}$.
The intrusion at $(i,k,\tau)$ is
\begin{equation}
\rho^{(i,k)}_\tau = \max\!\left(0,\ 1 - \sqrt{
\left(\frac{a^{(i,k)}_\tau}{m^{\mathrm{lon}}_i}\right)^{\!2} +
\left(\frac{\ell^{(i,k)}_\tau}{m^{\mathrm{lat}}_i}\right)^{\!2}}\;
\right) m^{\mathrm{lat}}_i,
\label{eq:score}
\end{equation}
and $u_t = \max_{i,k,\tau} \rho^{(i,k)}_\tau$. A single isotropic
clearance radius cannot represent both margins without either missing a
straight-ahead lead vehicle (a radius set by width alone) or
over-triggering on adjacent-lane traffic (a radius set by the vehicle's
half-diagonal); the decomposition in \eqref{eq:score} is required for
correctness, not merely convenient, a fact confirmed empirically during
development (Section~\ref{sec:setup}).

\subsection{Fallback manoeuvre}
\label{sec:mrm}
On triggering, the vehicle executes the minimum-risk manoeuvre defined
by UN Regulation No.~157 (Automated Lane Keeping Systems)
\S\,5.5.1~\cite{unr157}: decelerate inside the current lane with a
deceleration demand not exceeding $4.0\,\mathrm{m/s^2}$, lane-keeping
control otherwise unchanged, holding the vehicle stationary once
stopped. This is enforced by a proportional controller tracking the
target speed profile $v(t) = \max(v_0 - 4.0\,t,\,0)$ from the moment of
triggering.

%=======================================================================
\section{Experimental Setup}
\label{sec:setup}

\subsection{Simulation environment}
All rollouts are executed headless in MetaDrive, replaying scenarios
converted by ScenarioNet~\cite{scenarionet2023} from Argoverse~2 Motion
Forecasting and nuScenes. Non-ego traffic is controlled reactively
(rule-based, IDM), not replayed open-loop from the log: pure log
replay was found to roughly triple the no-fallback collision rate on a
1{,}000-scenario diagnostic sample, because logged agents cannot react
to an ego vehicle that deviates from its logged trajectory — precisely
the deviation a fallback manoeuvre produces. Reporting results under
non-reactive replay would therefore misrepresent the cost of stopping,
which is a central quantity this paper reports.

\subsection{Harm labelling}
\label{sec:harm}
A rollout is labelled harmful using the nuPlan closed-loop safety
definition~\cite{nuplan2024}: an at-fault collision, or a
time-to-collision below $0.95\,\mathrm{s}$ computed over a $3.0\,$s
horizon at $0.1\,$s resolution, not evaluated while the ego is
stationary and not evaluated against tracks behind the ego. Fault
attribution follows the nuPlan convention: a collision the ego suffers
while stationary, or while being struck from behind, is \emph{not}
at-fault and is never counted as a missed intervention; it is instead
recorded as an \emph{induced collision} — the concrete, measurable cost
of a false or premature trigger.

\subsection{Predictor}
The trajectory predictor is AutoBot, run online inside the closed loop
so that its input reflects the actual — possibly ego-deviated —
simulator history rather than the original log, which is necessary
once traffic is reactive. [Predictor training data, checkpoint
selection, and validation minADE to be finalised in the camera-ready
version alongside Section~\ref{sec:results}.]

\subsection{Splits and pre-registration}
Design decisions were made on a held-out development split, disjoint
from both the calibration split used to select $\hat\lambda$ and the
test split used only for the reported evaluation. The harm definition,
risk level $\alpha$, confidence level $\delta$, and the refutation
criteria for every hypothesis below were fixed and recorded before any
calibration- or test-split scenario was evaluated.

%=======================================================================
\section{Results}
\label{sec:results}

\emph{Results in this section are pending the full-scale calibration
and test-split evaluation ($n\!\approx\!2{,}000$ scenarios per split) and
are not yet reported. Diagnostic runs at small scale
($n\!\le\!300$) were used exclusively to validate the pipeline
described in Sections~\ref{sec:method}--\ref{sec:setup} — in particular,
to identify and correct two errors in the trigger-score geometry of
\eqref{eq:score} before any pre-registered evaluation — and are not
reported here as results.}

\subsection{H1: does open-loop ranking predict closed-loop ranking?}
[Pending.]

\subsection{H2: does calibrated triggering dominate hand-tuned thresholds?}
[Pending.]

\subsection{H3: does open-loop coverage collapse translate to closed-loop harm?}
[Pending.]

%=======================================================================
\section{Discussion}
\label{sec:discussion}
[To be completed alongside Section~\ref{sec:results}. Planned content
includes: sensitivity of $\hat\lambda$ to decision latency and decision
rate; the CVaR-style tail-severity extension motivated by
\cite{cvar2022}, replacing the mean missed-intervention rate in
\eqref{eq:guarantee} with a tail bound on a continuous
time-to-collision shortfall; and the limits of the reactive-IDM traffic
model as a proxy for real driver response to a decelerating vehicle.]

%=======================================================================
\section{Conclusion}
\label{sec:conclusion}
[To be completed once Section~\ref{sec:results} is populated.]

%=======================================================================
\section*{Acknowledgment}
The authors thank the maintainers of MetaDrive and ScenarioNet for
releasing the simulation infrastructure this work builds on.

%=======================================================================
\begin{thebibliography}{99}

\bibitem{vovk2005}
V.~Vovk, A.~Gammerman, and G.~Shafer, \emph{Algorithmic Learning in a
Random World}. Berlin, Germany: Springer-Verlag, 2005.

\bibitem{ltt2021}
A.~N.~Angelopoulos, S.~Bates, E.~J.~Cand\`es, M.~I.~Jordan, and L.~Lei,
``Learn then test: Calibrating predictive algorithms to achieve risk
control,'' \emph{arXiv preprint arXiv:2110.01052}, 2021.

\bibitem{crc2022}
A.~N.~Angelopoulos, S.~Bates, A.~Fisch, L.~Lei, and T.~Schuster,
``Conformal risk control,'' \emph{arXiv preprint arXiv:2208.02814},
2022.

\bibitem{robustcp2026}
K.~Rahaman, J.~V.~Deshmukh, A.~R.~Hota, and L.~Lindemann, ``When
environments shift: Safe planning with generative priors and robust
conformal prediction,'' in \emph{Proc. 8th Annu. Conf. Learn. Dyn.
Control (L4DC)}, 2026.

\bibitem{safepath2025}
A.~Doula, M.~M\"uhlh\"auser, and A.~S.~Guinea, ``SafePath: Conformal
prediction for safe LLM-based autonomous navigation,''
\emph{arXiv preprint arXiv:2505.09427}, 2025.

\bibitem{taskrelevant2022}
A.~Farid, S.~Veer, B.~Ivanovic, K.~Leung, and M.~Pavone, ``Task-relevant
failure detection for trajectory predictors in autonomous vehicles,''
in \emph{Proc. 6th Conf. Robot Learn. (CoRL)}, Auckland, New Zealand,
2022.

\bibitem{cvar2022}
M.~P.~Chapman, R.~Bonalli, K.~M.~Smith, I.~Yang, M.~Pavone, and
C.~J.~Tomlin, ``Risk-sensitive safety analysis using conditional
value-at-risk,'' \emph{IEEE Trans. Autom. Control}, vol.~67, no.~12,
pp.~6521--6536, Dec. 2022.

\bibitem{scenarionet2023}
Q.~Li, Z.~Peng, L.~Feng, Z.~Liu, C.~Duan, W.~Mo, and B.~Zhou,
``ScenarioNet: Open-source platform for large-scale traffic scenario
simulation and modeling,'' in \emph{Proc. 37th Conf. Neural Inf.
Process. Syst. (NeurIPS), Track Datasets and Benchmarks}, 2023.

\bibitem{nuplan2024}
N.~Karnchanachari \emph{et al.}, ``Towards learning-based planning: The
nuPlan benchmark for real-world autonomous driving,'' in \emph{Proc.
IEEE Int. Conf. Robot. Autom. (ICRA)}, 2024.

\bibitem{unr157}
United Nations Economic Commission for Europe, ``UN Regulation No.~157
--- Uniform provisions concerning the approval of vehicles with regard
to Automated Lane Keeping Systems,'' Rev.~1, 2025.

\end{thebibliography}

\end{document}
